\chapter{Adaptive Thresholding}
\label{ch:adaptive_thresholding}

\section{Motivation}
\label{sec:threshold_motivation}

In anomaly detection systems, the decision boundary that separates normal from 
abnormal behavior is often as important as the detector itself. In distributed 
microservice environments, however, workload characteristics are highly dynamic: 
traffic volume varies across services, diurnal patterns affect baseline latency, 
and deployments may temporarily shift normal operating conditions \cite{turn0search8}. 
Under such conditions, a fixed anomaly threshold often leads to either excessive false 
positives or missed incidents \cite{turn0search8}.

This issue is particularly important in operational monitoring, where frequent 
false alarms can lead to \emph{alert fatigue}, reducing trust in the detection 
system and delaying response to genuine incidents \cite{turn0search8}. Conversely, 
overly permissive thresholds may suppress early warning signals and allow 
performance degradation to propagate unnoticed. Therefore, the anomaly decision 
threshold should adapt to the recent behavioral context of each service.

Formally, let $e_t^{(i)}$ denote the anomaly score (e.g., reconstruction error) 
for service $S_i$ at time $t$, and let $\tau_t^{(i)}$ denote the corresponding 
decision threshold. An anomaly is declared when:

\[
y_t^{(i)} =
\begin{cases}
1, & e_t^{(i)} > \tau_t^{(i)} \\
0, & \text{otherwise}
\end{cases}
\]

The goal of adaptive thresholding is to learn or estimate $\tau_t^{(i)}$ such 
that detection remains robust under changing system conditions \cite{turn0search1}.

\section{Fixed vs Adaptive Thresholding}
\label{sec:threshold_fixed_vs_adaptive}

A \emph{fixed threshold} uses a constant decision boundary over time:

\[
\tau_t^{(i)} = \tau^{(i)}.
\]

This approach is simple and deterministic, but assumes that the statistical 
properties of the anomaly score remain stable. In practice, this assumption is 
often violated by burst traffic, seasonal load patterns, or gradual baseline 
drift \cite{turn0search8}.

By contrast, an \emph{adaptive threshold} evolves over time as a function of 
recent observations:

\[
\tau_t^{(i)} = f\big(e_{t-w:t}^{(i)}, z_t^{(i)}\big),
\]

where $e_{t-w:t}^{(i)}$ denotes a rolling history of anomaly scores and 
$z_t^{(i)}$ may include contextual variables such as load, latency, or Service 
Level Objective (SLO) status. Adaptive thresholding reduces false positives 
during legitimate spikes and improves sensitivity after sustained baseline 
shifts, at the cost of added algorithmic complexity \cite{turn0search8}.

For microservice monitoring, adaptive thresholding is preferable because service 
behavior is heterogeneous and non-stationary; a single static threshold rarely 
generalizes across all services and time periods \cite{turn0search8}.

\section{Reinforcement Learning for Threshold Adaptation}
\label{sec:threshold_rl}

\subsection{Reinforcement Learning Formulation}
\label{subsec:threshold_rl_formulation}

Adaptive thresholding can be formulated as a sequential decision problem, where 
an agent learns how to adjust the threshold in response to detection outcomes 
and service health signals.

For each service $S_i$, define the reinforcement learning tuple 
$\mathcal{M}^{(i)} = (\mathcal{X}, \mathcal{U}, P, R)$, where:

\begin{itemize}
    \item $\mathcal{X}$ is the state space, representing the current detection context;
    
    \item $\mathcal{U}$ is the action space, representing threshold adjustments;
    
    \item $P$ is the environment transition function;
    
    \item $R$ is the reward function.
\end{itemize}

A state may include recent detection quality and service health indicators:

\[
x_t^{(i)} =
\big[
\tau_t^{(i)},
\text{FPR}_t^{(i)},
\text{FNR}_t^{(i)},
\text{SLO}_t^{(i)}
\big],
\]

where $\text{FPR}$ and $\text{FNR}$ denote false positive and false negative 
rates, respectively, and $\text{SLO}_t^{(i)}$ reflects recent service-level 
objective violations.

The action $u_t^{(i)} \in \mathcal{U}$ modifies the threshold:

\[
\tau_{t+1}^{(i)} = \tau_t^{(i)} + u_t^{(i)}.
\]

The reward should penalize undesirable alerting behavior while encouraging 
accurate and operationally meaningful detection. A generic reward function is:

\[
R_t^{(i)} =
\alpha \cdot \mathbb{I}(\text{TP})
- \beta \cdot \mathbb{I}(\text{FP})
- \gamma \cdot \mathbb{I}(\text{FN})
- \delta \cdot \mathbb{I}(\text{SLO violation}),
\]

where $\alpha, \beta, \gamma, \delta > 0$ are tunable weights.

\subsection{Q-Learning}
\label{subsec:threshold_qlearning}

Q-learning is a model-free reinforcement learning algorithm that estimates the 
expected long-term utility of selecting action $u$ in state $x$ \cite{turn0search6}. 
The optimal action-value function is defined as:

\[
Q^*(x,u) = \mathbb{E}\left[\sum_{k=0}^{\infty} \gamma^k R_{t+k+1} \,\middle|\, x_t=x, u_t=u \right],
\]

where $\gamma \in (0,1)$ is the discount factor. The Q-values are updated using 
the Bellman equation:

\[
Q(x_t,u_t) \leftarrow Q(x_t,u_t) +
\eta \left[
R_{t+1} + \gamma \max_{u'} Q(x_{t+1},u') - Q(x_t,u_t)
\right],
\]

where $\eta$ is the learning rate.

Action selection typically follows an $\varepsilon$-greedy policy: with 
probability $\varepsilon$, the agent explores a random action, and with 
probability $1-\varepsilon$, it selects the action with the highest Q-value 
\cite{turn0search6}.

\subsection{Q-Learning Implementations and Variants}
\label{subsec:threshold_qlearning_variants}

Several variants of Q-learning exist depending on the complexity of the state 
space. \emph{Tabular Q-learning} is suitable when the state-action space can be 
discretized into a manageable number of buckets. \emph{Deep Q-Networks (DQN)} 
replace the Q-table with a neural approximator and are better suited for large 
or continuous state spaces. \emph{Double DQN} further reduces overestimation bias.

For adaptive thresholding in operational monitoring, tabular Q-learning is often 
sufficient and preferable. The decision process is low-dimensional, the learned 
policy remains interpretable, and the computational overhead is modest, which is 
particularly useful in lightweight or edge-oriented deployments \cite{turn0search1}.

\subsection{Disadvantages and Limitations}
\label{subsec:threshold_rl_disadvantages}

Despite its advantages, RL-based thresholding has practical limitations. First, 
the method suffers from a \emph{cold-start problem}: before sufficient feedback 
has been observed, the learned threshold policy may be unstable \cite{turn0search1}. 
Second, reward signals may be sparse, especially when true anomalies are rare. 
Third, the environment is non-stationary, meaning that a previously learned 
threshold policy may become suboptimal after architectural or workload changes.

These limitations imply that RL-based thresholding should be viewed as a policy 
adaptation mechanism rather than a fully autonomous replacement for operational 
judgment.

\section{Implementation}
\label{sec:threshold_implementation}

\subsection{Architecture}
\label{subsec:threshold_impl_architecture}

The adaptive thresholding system is implemented by the \texttt{ThresholdTuner}
class, which instantiates a tabular Q-learning agent for each monitored
service. The architecture comprises four layers.

\textbf{Q-Table.} The action-value function $Q(x,u)$ is stored as a nested
dictionary mapping \\ $(\text{service}, \text{state\_key}, \text{action\_index})
\mapsto \mathbb{R}$. The Q-learning parameters are:

\begin{itemize}
    \item Learning rate $\eta = 0.1$,
    \item Discount factor $\gamma = 0.95$,
    \item Initial exploration rate $\varepsilon = 0.1$ with multiplicative
          decay $0.995$ down to a minimum of $0.01$.
\end{itemize}

\textbf{State Encoder.} The continuous state vector $x_t^{(i)}$ is discretized
into a hashable key by bucketing each feature:

\begin{itemize}
    \item Current threshold $\tau_t^{(i)}$: rounded to one decimal place.
    \item False positive rate and false negative rate: buckets of width $0.05$.
    \item SLO violation rate (computed over the 20~most recent observations):
          buckets of width $0.05$.
\end{itemize}

This coarse discretization keeps the Q-table compact---typically a few hundred
entries per service---while capturing the operationally meaningful distinctions
between detection regimes.

\textbf{Action Space.} Five actions are available:
$\mathcal{U} = \{-0.5, -0.2, 0.0, +0.2, +0.5\}$, each representing a
threshold adjustment in the same units as the anomaly score. After applying an
action, the threshold is clamped to a minimum of~$0.1$ to prevent degenerate
behavior.

\textbf{Reward Function.} The reward $R_t^{(i)}$ is computed from three
components with configurable weights:

\begin{equation}
    R_t^{(i)} = \underbrace{(\alpha_p + \alpha_r) \cdot \mathbb{I}(\text{TP})}_{\text{correct detection}}
    \underbrace{- \beta \cdot \mathbb{I}(\text{FP})}_{\text{false positive}}
    \underbrace{- \gamma \cdot \mathbb{I}(\text{FN})}_{\text{missed anomaly}}
    \underbrace{- \delta \cdot \mathbb{I}(\text{SLO violation})+ \delta' \cdot \mathbb{I}(\text{SLO compliance})}_{\text{SLO term}},
\end{equation}

with defaults $\alpha_p = \alpha_r = 0.3$, $\beta = 1.0$, $\gamma = 2.0$,
$\delta = 3.0$, and $\delta' = 0.4$. The asymmetric penalties reflect the
operational reality that missed anomalies and SLO violations are more costly
than false positives.

\textbf{Persistence.} The Q-table, per-service thresholds, and confusion-matrix
counters (TP, FP, FN, TN) are maintained in memory during operation. For the
standard node profile (Chapter~\ref{ch:federated_learning}), this state can be
serialized to the local SQLite database, enabling crash recovery without
re-learning the threshold policy from scratch.

\subsection{Integration with the Detection Pipeline}
\label{subsec:threshold_impl_integration}

The \texttt{ThresholdTuner} operates as a closed-loop feedback controller
within the anomaly detection pipeline, as illustrated by the following
per-observation cycle.

\begin{enumerate}
    \item \textbf{Score Emission.} The LSTM-AE on an edge node computes the
    reconstruction error $e_t^{(i)}$ for the current observation window.

    \item \textbf{Threshold Comparison.} The error is compared against the
    current per-service threshold $\tau_t^{(i)}$ maintained by the
    \texttt{ThresholdTuner}. An anomaly is declared if
    $e_t^{(i)} > \tau_t^{(i)}$.

    \item \textbf{Feedback Ingestion.} The detection outcome---whether an
    anomaly was detected, whether the observation was truly anomalous (derived
    from SLO violation correlation via the \texttt{SLOTracker}), and whether an
    SLO breach occurred---is fed back to the tuner via
    \texttt{update\_feedback()}. This updates the internal confusion-matrix
    counters.

    \item \textbf{Threshold Adjustment.} At a configurable cadence (e.g.,
    every 10~observations), \texttt{tune\_threshold()} is invoked. The tuner
    computes the current state $x_t^{(i)}$, selects an action via the
    $\varepsilon$-greedy policy, updates the Q-table using the Bellman update,
    and applies the resulting threshold adjustment.
\end{enumerate}

This tight integration ensures that the threshold responds to recent detection
performance within minutes, not hours, while the decaying $\varepsilon$
gradually shifts the policy from exploration to exploitation as the agent
accumulates experience.

\subsection{Example: Dynamic CPU Utilization Alert Threshold}
\label{subsec:threshold_example}

Consider a service monitored over a 24-hour period (1\,440 observations)
using real-world metrics from the SMD dataset.

\textbf{Scenario~1: Static vs. EWMA thresholds.}
A static threshold is overly rigid. In our benchmarks
(Section~\ref{sec:eval_threshold}), a static threshold produced 72 false alerts
(FPR $\approx$ 5\%) and failed to detect genuine anomalies effectively
(F1 = 0.000). The EWMA baseline, which tracks the recent mean, performed worse
on this dataset: its threshold lagged behind normal metric variations,
generating 648 alerts (FPR $\approx$ 47\%).

\textbf{Scenario~2: Q-Learning exploration.}
The Q-learning agent begins with an initial threshold $\tau_0 = 0.1$. Over the
first 24~hours, it observes that its high threshold suppresses all alerts
(FPR = 0.000), which yields positive reward for correctly identifying the
abundant normal traffic (true negatives). However, it also completely misses
the injected anomalies (F1 = 0.000). The agent accumulates positive reward
from true negatives (reaching a cumulative reward of 552.4), but the reward
signal is dominated by the majority normal class.

\textbf{Key observation.} The Q-learning agent's $\varepsilon$ parameter successfully
decays from 0.10 to 0.05, establishing the expected exploration-exploitation
shift. However, the empirical results show that the reward function must heavily
penalize false negatives to force the threshold downwards when true negatives
dominate the dataset. Benchmark~E (Section~\ref{sec:eval_threshold}) records
the full 24-hour trajectory, highlighting the importance of extended
training horizons and asymmetric reward tuning for adaptive thresholding
in highly imbalanced real-world telemetry.
